\section{Числовые характеристики случайных величин.}
Пускай у нас всюду в этой секции задано некоторое вероятностное пространство $(\Omega, \mathcal{F}, \P)$.
\subsection{Математическое ожидание как интеграл Лебега.}
\subsubsection{Основные свойства.}
\begin{definition}
    Пусть $\xi: \Omega \to \R$ -- случайная величина.
    Будем называть математическим ожиданием случайной величины (и обозначать $\Ex \xi$) интеграл Лебега по $\P$ на $\Omega$:
    \[
        \Ex \xi = \int_{\Omega} \xi d\P.
    \]
    (естественно, математического ожидания может не существовать).
\end{definition}
Назовём несколько основных свойств математического ожидания:
\begin{enumerate}
    \item $\Ex c = c, \Ex_{I_A} = \P(A)$.
    \item $\Ex \alpha \xi + \beta \eta = \alpha \Ex \xi + \beta \Ex \eta$.
    \item Если две случайные величины $\xi$ и $\eta$ совпадают почти всюду, то $\Ex \xi = \Ex \eta$.
    \item Если $\xi \leq \eta$ почти всюду, то и $\Ex \xi \leq \Ex \eta$.
    \item Если $\xi \geq 0$ почти всюду, но $\Ex \xi = 0$, то $\xi = 0$ почти всюду.
    \item Пусть $\eta = \phi(\xi)$, где $\xi$ -- случайная величина, а $\phi$ -- борелевская функция.
    Тогда $\Ex \eta = \int_{\Omega} \phi(\xi)d\P = \Ex \phi(\xi)$.
    \item $|\Ex \xi| \leq \Ex |\xi|$.
\end{enumerate}
\subsubsection{Неравенства Йенсена и Ляпунова.}
\begin{theorem}[Неравенство Минковского]
    Для случайных величин $\xi$ и $\eta$ и некоторого $r > 1$ верно следующее:
    \[
        \|f + g\|_{r} \leq \|f\|_{r} + \|g\|_{r}.
    \]
\end{theorem}
\begin{theorem}[Неравенство Гёльдера]
    Для случайных величин $f \in L_p$ и $g \in L_{p'}$ справедливо следующее неравенство:
    \[
        \|f \cdot g\|_{1} \leq \|f\|_{p} \cdot \|g\|_{p'}.
    \]
\end{theorem}
\begin{theorem}[Неравенство Ляпунова]
    Для случайной величины $\xi$ и $u > v > 0$ справедливо следующее:
    \[
        \|\xi\|_{u} \geq \|\xi\|_{v}.
    \]
\end{theorem}
\begin{theorem}[Неравенство Йенсена]
    Для выпуклой вниз борелевской функции $g$ и случайной величины $\xi$ справедливо следующее неравенство:
    \[
        \Ex g(\xi) \geq g\bigr(\Ex \xi \bigr).
    \]
\end{theorem}
\subsubsection{Математическое ожидание независимых и некоррелированых случайных величин.}
\begin{theorem}
    Пусть $\xi$ и $\eta$ -- независимые случайные величины.
    Тогда, в случае существования у них первого момента, справедливо следующее равенство:
    \[
        \Ex \xi \eta = \Ex \xi \Ex \eta.
    \]
\end{theorem}
\begin{proof}
    Для простых функций данное утверждение очевидно.
    И действительно, если $\xi = \sum\limits_{k = 1}^{n} x_k I_{A_k}$ и $\eta = \sum\limits_{r = 1}^{m} y_k I_{B_r}$, то
    \[
        \Ex \xi \eta = \sum\limits_{k, r = 1, 1}^{n, m} x_k y_r \P(A_k B_r) = \sum\limits_{k, r = 1, 1}^{n, m} x_k y_r \P(A_k) \P(B_r) = \sum\limits_{k = 1}^n x_k \P(A_k) \cdot \sum\limits_{r = 1}^m y_k \P(B_r) = \Ex \xi \cdot \Ex \eta.
    \]
    В общем случае, для каждой неотрицательной случайной величины существуют последовательности $\xi_n$ и $\eta_n$, которые поточечно сходятся к $\xi$ и $\eta$ (без ограничения общности можно считать сходимость монотонной).
    Тогда, $\Ex \xi_n \eta_n =  \Ex \xi_n \Ex \eta_n$ -- по предыдущему пункту, и, совершая предельный переход по теореме Леви, получаем утверждение для неотрицательных функций. \\
    В общем случае для интегрируемых функций можно воспользоваться теоремой Лебега о мажорируемой сходимости и получить тот же результат.
\end{proof}
\subsubsection{Экспоненциальное неравенство Маркова, неравенство Чернова.}

\subsubsection{Неравенство Чебышёва для $i.i.d.$-последовательности.}

\subsubsection{Функции моментов, дисперсия, ковариация.}
\begin{definition}
    Для случайных величин $\xi$ и $\eta$ с конечным вторым моментом можно определить скалярное произведение следующим образом:
    \[
        (\xi, \eta) = \Ex \xi \eta.
    \]
\end{definition}
\begin{reminder}
    Пространство $L_2(\P)$ является гильбертовым пространство со скалярным произведением, определённым выше.
\end{reminder}
\begin{definition}
    Пусть $\xi$ -- случайная величина.
    Будем обозначать $\overset{\circ}{\xi}$ центрированную $\xi$, то есть $\overset{\circ}{\xi} = \xi - \Ex \xi$.
\end{definition}
\begin{definition}
    Будем называть $k$-м моментов математическое ожидание $\xi^k$, то есть $\Ex \xi^k$.
\end{definition}
\begin{definition}
    Будем называть $k$-м центральным моментом величину $\Ex \overset{\circ}{\xi} ^k$.
\end{definition}
\begin{definition}
    Будем называть $k$-м абсолютным моментом величину $\Ex |\xi|^k$
\end{definition}
\subsubsection{Корреляция и её свойства.}
\subsection{Математическое ожидание как интеграл Лебега-Стилтьеса.}
\subsubsection{Математическое ожидание дискретных и непрерывных случайных величин, примеры.}
\subsubsection{Маргинальные распределения.}
\subsection{Условное математическое ожидание.}
\subsubsection{Достаточное условие существования и единственности условного математического ожидания.}
\subsubsection{Теорема о наилучшем квадратичном прогнозе.}
\section{Преобразования случайных величин.}
\subsection{Характеристические функции.}
\begin{definition}
    Характеристической функцией случайной величины $\xi$ будет называть преобразование Фурье меры $\P_\xi$, то есть
    \[
        \phi_{\xi}(t) = \Ex e^{it \xi}.
    \]
\end{definition}
Таким образом, для абсолютно непрерывных случайных величин характеристические функции имеют вид
\[
    \phi_{\xi}(t) = \int_{\R} f_{\xi}(x)e^{itx}dx.
\]
А для дискретных случайных величин характеристические функции имеют вид
\[
    \phi_{\xi}(t) = \sum\limits_{k} \P(\xi = y_k)\cdot e^{ity_k}.
\]
Характеристические функции обладают следующим набором свойств:
\begin{enumerate}
    \item $|\phi(t)| \leq 1$. И действительно: $|\phi(t)| = |\Ex e^{it \xi}| \leq \Ex |e^{it \xi}| = \Ex 1 = 1$.
    \item Характеристическая функция является равномерно непрерывной на всей числовой прямой: $\forall h \in \R \hookrightarrow |\phi_\xi(t + h) - \phi_\xi(t)| \leq \Ex |e^{i(t + h)\xi} - e^{it \xi}| = \Ex |e^{ih} - 1| \ra 0, h \ra 0$ по теореме Лебега о мажорируемой сходимости.
    \item Если $\eta = a \xi + b$, то $\phi_{\eta}(t) = \Ex e^{it \eta} = \Ex e^{it (a \xi + b)} =e^{itb} \Ex e^{i (ta)} = e^{itb}\phi_\xi (ta) $.
    \item $\phi^{(k)}_\xi(t)\big|_{t = 0} = i^k \int x^k e^{it \xi} dF_\xi \big|_{t = 0} = \Ex \xi^k$
    \item Если $\xi$ и $\eta$ -- независимые случайные величины, то и функции $e^{it\xi}$ и $e^{it \eta}$ -- независимы, а следовательно $\phi_{\xi + \eta}(t) = \Ex e^{it (\xi + \eta)} = \Ex e^{it\xi}\cdot \Ex e^{it \eta} = \phi_\xi(t)\phi_\eta(t)$.
\end{enumerate}
\begin{example}
    Посчитаем характеристические функции некоторых известных распределений:
    \begin{itemize}
        \item Для $\xi \sim Be(p)$ очевидно, что $\phi_\xi(t) = q + pe^{it}$.
        \item Для $\xi \sim Binom(n, p)$ в силу того, что $\xi = \sum\limits_{k = 1}^n y_k$ -- i.i.d. распределённых $Be(p)$, то $\phi_\xi(t) = (p + qe^{it})^n$.
        \item Для $\xi \sim N(0, 1)$:
        \[
            \phi_{\xi}(t) = \int \dfrac{1}{\sqrt {2\pi}} e^{-x^2/2} e^{itx}dx.
        \]
        Так как $\phi'(t) = \dfrac{-i}{\sqrt {2\pi}} \int e^{itx}de^{-x^2/2} = \dfrac{i}{\sqrt {2\pi}} \int e^{-x^2/2}de^{itx} = -t \phi(t)$.
        Очевидно, что $\phi(t) = Ce^{-t^2/2}$. Из нормировки в нуле $\phi(0) = 1$ мы получаем искомую функцию:
        \[
            \phi_{xi}(t) = e^{-t^2/2}.
        \]
        \item Таким образом, если $\xi \sim N(m, \sigma^2) = m + \sigma N(0, 1)$, то $\phi_\xi(t) = e^{itm - \dfrac{(\sigma t)^2}{2}}$.
    \end{itemize}
\end{example}
\begin{note}
    Также можно заметить, что сумма независимых нормальных распределений также распределена нормально.
\end{note}
% Формулу обращения и заебись.
\begin{theorem}[О единственности]
    Для случайных величин $\xi$ и $\eta$ следующие условия эквивалентны:
    \begin{enumerate}
        \item $\phi_\xi = \phi_\eta$.
        \item $F_{\xi} = F_{\eta}$.
    \end{enumerate}
\end{theorem}
\begin{proof}
    Надо-надо... Наверное?
\end{proof}
\begin{note}
    Слабая сходимость распределений эквивалентна поточечной сходимости их характеристических функций.
\end{note}
\begin{note}
    Последовательность характеристических функций, в общем случае, может сходится не к характеристической функции -- например $\phi_{U(-n, n)} \ra \phi_{I_{\{0\}}}$.
\end{note}
\subsection{Производящие функции.}
\begin{definition}
    Пусть $\xi$ -- натурально-значная случайная величина (то есть в $\N_0$).
    Тогда производящей функцией $\xi$ будем называть
    \[
        G_{\xi}(z) = \Ex z^{\xi} = \sum\limits_{k = 0}^{+\infty} \P(\xi = k)z^k.
    \]
\end{definition}
Производящая функция случайной величины имеет следующие свойства:
\begin{enumerate}
    \item Ряд абсолютно сходится при $|z| < 1$, поскольку $G(1) = 1$ (это просто мера вероятностного пространства).
    \item Производящая функция однозначно определяет распределение: $\P(\xi = k)\cdot k! = G^{(k)}(0)$.
    \item $G^{m}(1) = \Ex (\xi \cdot \ldots (\xi - m + 1))$.
    \item $\Ex \xi = G'(1)$ и $\D \xi = G''(1) + G'(1) - G'(1)^2$.
    \item $G_{a\xi + b}(z) = \Ex z^{a\xi + b} = z^b G_{\xi}(z^a)$.
    \item Если $\xi$ и $\eta$ -- независимые случайные величины, то $G_{\xi + \eta} = G_{\xi}G_{\eta}$.
\end{enumerate}
\begin{example}
    Приведём несколько примеров производящих функций:
    \begin{itemize}
        \item Для распределения Бернулли с параметром $p$ справедливо $G_{Be(p)} = q + pz$.
        \item Для распределения Пуассона c параметром $\lambda$:
        \[
            G(z) = \sum\limits_{k = 0}^{+\infty} z^k \cdot \dfrac{\lambda^k e^{-\lambda}}{k!} = e^{-\lambda} \sum\limits_{k = 0}^{+\infty} \dfrac{(\lambda z)^k}{k!} = e^{\lambda(z - 1)}.
        \]
    \end{itemize}
\end{example}
\subsection{Характеристические функции моментов.}
\begin{definition}
    Производящей функцией моментов для случайной величины $\xi$ будем называть
    \[
        H_{\xi}(t) = \Ex e^{t\xi}.
    \]
\end{definition}
Она обладает следующими базовыми свойствами:
\begin{enumerate}
    \item Она однозначно определяет распределение $\xi$.
    \item $H_{\xi}(0) = 0, h_{\xi}(t) > 0$ при всех $t$, где эта функция определена.
    \item $H_{a\xi + b}(t) = e^{bt}H_{xi}(at)$.
    \item Все моменты $\Ex \xi^k$ определены, а по ним однозначно восстанавливается распределение.
    \item $\Ex \xi^k = H^{(k)}_\xi(0)$.
    \item Аналогично, для независимых случайных величин характеристическая функция моментов суммы равна произведению характеристических функций моментов элементов суммы.
\end{enumerate}
\begin{example}
    Приведём несколько примеров характеристических функций моментов.
    \begin{itemize}
        \item Для $\xi \sim N(0, 1)$ она равна $H_\xi(t) = e^{t^2/2}$.
        \item Для гамма-распределения $\Gamma(\alpha, \lambda)$ она равна $H_{\xi}(t) = (\lambda/\lambda - t)^\alpha$.
        Пусть $X \sim \Gamma(\alpha, \lambda)$. Тогда у $X$ есть плотность $f_X (x) = \dfrac{\lambda^\alpha}{\Gamma(\alpha)}x^{\alpha - 1}e^{-\lambda x} \cdot I_{x > 0}$.
        Характеристическая функция моментов:
        \[
            H_X(t) = \int_0^{+\infty} e^{tx} \cdot \dfrac{\lambda^\alpha}{\Gamma(\alpha)} x^{\alpha - 1}e^{-\lambda x} dx = \dfrac{\lambda^\alpha}{(\lambda - t)^\alpha} \int_0^{\infty} \dfrac{(\lambda - t)^\alpha}{\Gamma(\alpha)}x^{\alpha - 1} e^{-x(\lambda - t)}dx.
        \]
        Отсюда видно, что характеристическая функция моментов для гамма-распределения существует при $\lambda > t$.
        Поскольку это, по сути, интеграл от плотности по всему пространству, то
        \[
            H_X(t) = \dfrac{\lambda^\alpha}{(\lambda - t)^\alpha}.
        \]
    \end{itemize}
    \begin{note}
        Из этого можно найти характеристическую функцию гамма-распределения $\phi_X(t) = \dfrac{\lambda^\alpha}{(\lambda - it)^\alpha}$, так как $\phi_X(t) = H_X(it)$ (формально).
    \end{note}
\end{example}
\subsection{Интегральная теорема Муавра-Лапласа.}
\begin{theorem}[Муавр, Лаплас, интегральная версия]
    Пусть $\mu_n \sim Binom(n, p)$.
    Тогда $\Ex \mu_n = np, \D \mu_n = npq$.
    Рассмотрим случайную величину $\eta_n = \dfrac{\mu_n - \Ex \mu_n}{\sqrt {\D \mu_n}} = \dfrac{\mu_n - np}{\sqrt {npq}}$.
    Для функции распределения $\eta_n$ справедливо следующее утверждение: \[\P(\eta_n \leq x) F_{\eta_n}(x) \rr F_{N(0, 1)}(x) = \dfrac{1}{\sqrt {2\pi}}\int_{-\infty}^x e^{-t^2/2}dt.\]
\end{theorem}
\begin{proof}
    По сути, утверждение теоремы состоит в $\eta_n \xrightarrow{d} N(0, 1), n \ra +\infty$.
    По теореме Леви о непрерывности, это эквивалентно поточечной сходимости характеристических функций всюду на $\R$.
    Рассмотрим характеристическую функцию $\eta_n$:
    \begin{multline*}
        \phi_{\eta_n}(t) = \phi_{\frac{\mu_n}{\sqrt {npq}} - \frac{np}{\sqrt {npq}}}(t) = e^{-\frac{itnp}{\sqrt {npq}}}\cdot \phi_{\mu_n}\big(\frac{t}{\sqrt {npq}}\big) = e^{-\frac{itnp}{\sqrt {npq}}}(q + pe^{\frac{it}{\sqrt {npq}}})^n
        = (qe^{-\frac{itp}{\sqrt {npq}}} + pe^{\frac{it(1 - p)}{\sqrt {npq}}})^n = \\ = (q(1 - itp/\sqrt {npq} - (pt)^2/2npq) + p(1 - qt/\sqrt {npq} - (qt)^2/2npq) + o(1/n))^n = \\ =
        (1 - t^2/2n + o(1/n))^n \ra e^{-t^2/2}, n \ra +\infty.
    \end{multline*}
    Что и доказывает утверждение теоремы.
    Сходимость равномерная, поскольку функция распределения нормальной случайной величины -- непрерывна.
\end{proof}
\section{Нормальный случайный вектор.}
\begin{definition}
    Пусть есть случайный вектор $\vec{\xi}$, его математическое ожидание равно $\vec{m} = \Ex \vec{\xi}$. \\
    Рассмотрим матрицу ковариации этого случайного вектора, которая определяется как $R = \{\cov \xi_i, \xi_j\}_{i, j = 1, 1}^{n, n}$.
    Будем говорить, что $\vec{\xi}$ -- нормальный случайный вектор с параметрами $\vec{m}$ и $R$, если $\det R \neq 0$ и он имеет плотность:
    \[
        p_{\vec{\xi}}(\vec{x}) = \dfrac{1}{(\sqrt {2\pi})^n \sqrt {|\det R|}} \exp\biggr( -\dfrac{1}{2}(\vec{x} - \vec{m})^T R^{-1}(\vec{x} - \vec{m}) \biggr).
    \]
\end{definition}
\begin{lemma}
    Заметим, что для нормального случайного вектора некоррелированность и независимость компонент эквивалентна -- совместная плотность просто представляется в виде произведения плотностей.
\end{lemma}
\begin{note}
    Заметим, что если $\xi \sim N(\vec{m}, R)$ и $\eta = B \xi$, то плотность преобразуется таким образом, что $\eta \sim (B \vec{m}, BRB^T)$:
    Это так, поскольку $\xi = B^{-1}\eta$ и \[|J| = \dfrac{1}{|\det B|} = \dfrac{1}{\sqrt {\det B \det B^T}}.\]
    Теперь прямо посмотрим на плотность $\eta$:
    \[
        p_{\vec{\eta}}(\vec{x}) = \dfrac{1}{(\sqrt{2\pi})^{n} \sqrt{\det|BRB^T|}} \exp (-\dfrac{1}{2}(B^{-1}\vec{x} - \vec{m})^TR^{-1}(B^{-1}\vec{x} - \vec{m})).
    \]
    Выражение в показателе экспоненты преобразуется как
    \[
        (B^{-1}(\vec{x} - B\vec{m}))^T R^{-1} B^{-1}(\vec{x} - B\vec{m}) = (\vec{x} - B\vec{m})(BRB^T)^{-1}(\vec{x} - B\vec{m}).
    \]
    что и показывает справедливость нашего утверждения.
\end{note}
\subsection{Характеристическая функция многомерного гауссовского распределения.}
\begin{definition}
    Модифицируем определение характеристической функции на случайные вектора.
    Собственно, для этого не надо выдумывать что-то сверхестественное -- просто используем многомерное преобразование Фурье:
    \[
        \phi_{\vec{\xi}}(\vec{t}) = \Ex e^{i \langle \vec{t}, \vec{\xi} \rangle}.
    \]
\end{definition}
\begin{lemma}
    Характеристическая функция невырожденного гауссовского случайного вектора $\vec{\xi} \sim N(\vec{m}, R)$ имеет вид
    \[
        \phi_{\vec{\xi}}(\vec{t}) = e^{i \langle \vec{m}, \vec{t} \rangle - \frac{1}{2} \vec{t}^T R \vec{t}}.
    \]
\end{lemma}
\begin{proof}
    В силу положительной определённости матрицы $R$ сущесвтует такая ортогональная $U$, что $URU^T = diag(\sigma_1^2, \ldots \sigma_n^2)$.
    Рассмотрим тогда случайный вектор $\vec{\eta} = U\vec{\xi}$.
    Он распределён как $N(U\vec{m}, diag(\sigma_1^2, \ldots, \sigma_n^2))$.
    Его компоненты -- некоррелированные случайные величины, а следовательно -- независимые.
    Тогда его харфункция является произведение харфункций его компонент и равна
    \[
        \phi_{\vec{\eta}}(\vec{t}) = e^{i \langle U\vec{m}, \vec{t} \rangle- \frac{1}{2} \vec{t}^T diag(\sigma_1^2, \ldots) \vec{t}}.
    \]
    Поскольку мы знаем как меняются харфункции при линейных преобразованиях случайных величин (а, следовательно, и случайных векторов -- это несложно показать по определению), обратное пребразование $\vec{\xi} = U^{-1}\vec{\eta}$ даст нам утверждение леммы.
\end{proof}
\begin{note}
    Характеристические функции случайных векторов во многом имеют свойства, аналогичные свойствам характеристтических функций случайных величин.
\end{note}
\subsection{Обобщённый нормальный вектор. Линейное преобразование обобщённого нормального вектора.}
\begin{definition}
    Требование невырожденности матрицы ковариации может быть слегка излишним, поэтому мы можем сделать такое обобщение, и называть $\vec{\xi}$ обобщённым нормальным вектором, если
    \[
        \phi_{\vec{\xi}}(\vec{t}) = e^{i \langle \vec{t}, \vec{m} \rangle - \frac{1}{2}\vec{t}^T R \vec{t}}.
    \]
    для некоторого вектора $\vec{m}$ и неотрицательно определённой симметрической матрицы $R$.
\end{definition}
\begin{lemma}
    Пусть $\vec{\xi} \sim (\vec{m}, R)$ -- обобщённый гауссовский вектор. Тогда если $\vec{\eta} = A\vec{\xi} + \vec{b}$, то
    $\vec{\eta}$ -- также обобщённый нормальный вектор, который распределён с параметрами $N(A\vec{m} + \vec{b}, ARA^T)$.
\end{lemma}
\begin{proof}
    \[
        \phi_{\vec{\eta}}(\vec{t}) = \Ex e^{i \langle \vec{t}, A\vec{\xi} + \vec{b} \rangle} = e^{i\langle \vec{t}, \vec{b} \rangle} \Ex e^{i \vec{t}^T A \vec{\xi}} = e^{i\langle \vec{t}, \vec{b} \rangle} \Ex e^{i (A^T\vec{t})^T \vec{\xi}}
        = e^{i \langle \vec{t}, \vec{b} \rangle} \phi_{\vec{\xi}}(A^T \vec{t}) = e^{i \langle \vec{t}, \vec{b} \rangle} \cdot e^{i \langle \vec{t}, A\vec{m} \rangle} \cdot e^{-\frac{1}{2} \vec{t}^T ARA^T \vec{t}}.
    \]
\end{proof}